\titledquestion{布朗运动的等价定义}
\begin{parts}
    % @QUESTION
    \part 证明过程 $(X_t)_{t\in \R^+}$ 为布朗运动（Brownian motion）当且仅当：
        \begin{itemize}
        \item[i')] $X_0=0$,
        \item[ii')] $(X_t)_{t\in \R^+}$ 为零均值高斯过程（centered Gaussian process）,
        \item[iii')] ${\rm Cov}(X_s,X_t)=\inf(s,t).$
        \end{itemize}
    % @QUESTION_END
    \begin{xsolution}
    % @SET_START
    首先证明必要性方向：布朗运动 $\Rightarrow$ 高斯过程。
    % @LIST_START
    \begin{enumerate}
        \item 增量的独立性结构：  
           % @ATOM
           % @PRECOND
           对于 \(s \leq t\)，将 $X_t$ 分解为 \(X_t = X_s + (X_t - X_s)\)。
           % @ARGUMENT
           由 \(X_s\) 与增量 \(X_t - X_s\) 的独立性，
           % @ARGUMENT_END
           协方差可化简为：  
           \[
           \mathrm{Cov}(X_s, X_t) = \mathrm{Cov}(X_s, X_s) + \mathrm{Cov}(X_s, X_t - X_s) = \mathrm{Var}(X_s) + 0 = s
           \]  
           直接得出 
           % @OUTCOME
           \(\mathrm{Cov}(X_s, X_t) = \min(s,t)\)。
           % @ATOM_END
        \item 高斯性质：  
           % @ATOM
           % @PRECOND
           标准布朗运动具有高斯增量（Gaussian increments）且 
           % @PRECOND
           \(X_0 = 0\)，
           % @PRECOND_END
           因此任意有限维线性组合 \((X_{t_1}, \dots, X_{t_n})\) 均可表示为
           % @ARGUMENT
           独立高斯增量的线性组合，故其服从高斯分布。  
           % @OUTCOME
           因此该过程为零均值高斯过程。
           % @ATOM_END
    \end{enumerate}    
    % These two points essentially exploit \emph{independence of increments} and \emph{normality} of increments.
    % @LIST_END
    其次证明充分性方向：
    % @LIST_START
    \begin{enumerate}
    \item 增量协方差计算：  
       % @ATOM
       % @PRECOND
       利用性质 ${\rm Cov}(X_s,X_t)=\inf(s,t)$，
       % @PRECOND
       对于 \(s < t \leq u < v\)，计算协方差 $\text{Cov}(X_t - X_s, X_v - X_u)$：
       % @ARGUMENT:CALCUL
       \[
       \text{Cov}(X_t - X_s, X_v - X_u) = \min(t,v) - \min(t,u) - \min(s,v) + \min(s,u)
       \]
       此式恒等于 0，证明了
       % @OUTCOME
       不相交时间段的增量是不不相关的（uncorrelated），进而由于是高斯向量（Gaussian vector）可知其相互独立。
       % @OUTCOME_END
       % @ATOM_END
    \item 增量方差计算：  
       % @ATOM
       % @PRECOND
       计算增量 $X_t - X_s$ 的方差：
       % @PRECOND_END
       对于 \(s < t\)：
       \[
       \text{Var}(X_t - X_s) = t + s - 2s = t - s
       \]
       % @ARGUMENT
       因此增量的分布仅依赖于时间差 \(t-s\)。
       % @OUTCOME
       这表明增量具有平稳性（stationary），
       % @OUTCOME
       且服从分布 \(\mathcal{N}(0, t-s)\)。
       % @ATOM_END
    结合上述两点：
    % @SET_START
    % @ATOM
    % @PRECOND
    不相交增量相互独立，且
    % @PRECOND
    增量服从方差为 $t-s$ 的零均值高斯分布，即 \(\mathcal{N}(0, t-s)\)，
    % @ARGUMENT
    加上轨道连续性（continuity of paths）假设，
    % @ARGUMENT_END
    即可得到
    % @OUTCOME 
    布朗运动的经典定义。
    % @ATOM
    % @PRECOND
    不相交增量相互独立，且
    % @PRECOND
    增量服从方差为 $t-s$ 的零均值高斯分布，即 \(\mathcal{N}(0, t-s)\)，
    % @ARGUMENT
    通过 Kolmogorov 修正定理（Kolmogorov's continuity theorem）取其连续版本，
    % @ARGUMENT_END
    亦可导出
    % @OUTCOME 
    布朗运动的经典定义。
    % @SET_END
    \end{enumerate}
    % @LIST_END
    % @SET_END
    \end{xsolution}

    % @QUESTION
    \part 证明：若 $W_t$ 为布朗运动，则 $ \{X_t=t\;W_{\frac
    1t}, t>0, X_0=0 \}$
    亦为布朗运动（时间反转/时间反转变换 Time Inversion）。
    \begin{xsolution}
    % @SET_START
    % @ATOM
    对于在原点 0 处的连续性：
    % @PRECOND
    计算 $\lim_{t \to 0^+} t W_{1/t}$：
    % @PRECOND_END
    \[
    \lim_{t \to 0^+} t W_{1/t} = \lim_{s \to \infty} \frac{W_s}{s} = 0 \quad \text{a.s. (几乎必然)}
    \]
    % @ARGUMENT
    根据布朗运动的强大数定律（Law of large numbers）：
    % @OUTCOME 
    $\lim_{t \to 0^+} t W_{1/t} = 0$ a.s.
    % @ATOM
    \\
    对于零均值高斯性质：  
    % @PRECOND
    对任意给定时刻 \(t_1, \dots, t_n\)，随机向量 \((X_{t_1}, \dots, X_{t_n})\) 
    % @ARGUMENT
    均为 \(W_{1/t_i}\) 的线性组合，
    % @OUTCOME 
    故服从零均值高斯分布。
    \\
    % @ATOM
    对于协方差结构：  
    % @ARGUMENT:CALCUL
    \[
    \mathbb{E}[X_t X_s] = ts \cdot \min\left(\frac{1}{t}, \frac{1}{s}\right) = \min(t,s),
    \]
    % @OUTCOME 
    即 $\mathbb{E}[X_t X_s] = \min(t,s)$。
    \\
    % @ATOM
    综上所述：
    % @PRECOND
    轨道连续性：$\lim_{t \to 0^+} X_t = 0$ a.s.，
    % @PRECOND
    高斯性：对任意时刻集合，向量 \((X_{t_1}, \dots, X_{t_n})\) 均服从零均值高斯分布，
    % @PRECOND
    协方差结构：$\mathbb{E}[X_t X_s] = \min(t,s)$，
    % @PRECOND
    结合初始条件 \(X_0 = 0\)，
    % @ARGUMENT
    由布朗运动的等价判定性质（part a 结论）
    % @OUTCOME
    可知 \(X\) 亦为布朗运动。
    % @SET_END
    \end{xsolution}
\end{parts}